\label{sec:methods}

\subsection{Simulation architecture}
Numeric propagation, camera sensing, cloud occlusion, and reward computation live in \texttt{backend/simulation/}.
Rendering (\texttt{backend/render/}) is view-only and consumes \texttt{SimulationStateSeries} from \texttt{run\_simulation()}---no duplicate physics in export paths.

The current mission is a single-satellite 2D polar overflight with a primary nadir camera and a secondary wide-field camera.
Altitude is sampled per episode; targets lie on a meridian stripe (see \texttt{environment\_definition/constants/MISSION.py}).

\subsection{Observation and action}
Attitude dynamics are integrated in \texttt{SimulationStepper} with step $\Delta t = \SI{0.4}{\second}$ (\texttt{SIMULATION.simulation\_timestep}).
The policy outputs three-axis reaction-wheel torque commands bounded by plant limits.
Observations are assembled via explicit feature selection (\texttt{controller\_observation.py}); dimensions depend on enabled vision keys.

\subsection{Reward}
Reward terms combine distance-to-target shaping, area coverage and novelty, image-quality capture at discrete shutter events, and wheel-dynamics penalties.
Weights and gates are defined in \texttt{autonomous\_control/reward.py} and \texttt{docs/presentation/machine-learning.md}.
Reported numbers must match the code revision cited in the appendix.

\subsection{Learning algorithm}
Training uses MPO (\texttt{MPOAgent} in \texttt{controller\_agent.py}) with episode rollouts from \texttt{EpisodeRunner}.
Warmup and evaluation share the same simulation configuration as training.
Hyperparameters are recorded in \texttt{mpo\_config.py} and the presentation slides under \texttt{docs/presentation/}.

\subsection{Software architecture}
The codebase separates numeric simulation, learning, and visualization so train, eval, and export paths cannot silently diverge.
\begin{itemize}
  \item \textbf{Canonical episodes.} \texttt{backend/simulation/run\_simulation.py} is the single source of truth for episode-level world state; train, eval, and render consume \texttt{SimulationStateSeries} rather than re-running parallel integrators.
  \item \textbf{Timestep stack.} Attitude integration uses \texttt{SIMULATION.simulation\_timestep}; controller updates may differ from the physics step (nearest integer multiple, surfaced at runtime). Episode horizon scales with orbit window and \texttt{sat\_motion\_span\_scale}.
  \item \textbf{Seeded randomness.} Per-workflow \texttt{seed} derives deterministic altitude and cloud layouts (\texttt{derive\_seed}); warmup episodes use per-index derived seeds. Warmup bundles are fingerprinted and cached so repeated training runs skip redundant rollouts when the mission fingerprint matches.
  \item \textbf{Notebook cycle.} S01 notebooks~01--08 prototype features; stable logic is promoted to \texttt{s01\_utils/} (e.g.\ \texttt{training\_workflow.py}, baseline overflight). Run artifacts land under \texttt{autonomous\_control/models/} with config snapshots, telemetry, and video exports.
  \item \textbf{Parallelization.} \texttt{EpisodeRunner} supports serial rollouts for debugging and reproducibility; parallel batching is a throughput trade-off deferred where it would complicate artifact parity.
\end{itemize}
Constants and hyperparameters are maintained in \texttt{docs/presentation/} as the live record aligned with code.

\subsection{Baselines and evaluation}
Compare learned policies against documented baselines (e.g.\ zero-torque coast and scripted overflight sweeps in notebook~07).
Metrics should include capture count, mean image quality, cloud-blocked fraction, and attitude safety violations---extracted from run artifacts (\texttt{run\_log.md}, \texttt{summary\_metrics.json}) rather than re-simulated ad hoc.
Matched comparisons require the same altitude draw, cloud formation, and target layout per episode; the shared simulation kernel enforces this contract.
